% Generated draft: requires mathematical and source review.
\documentclass[11pt]{article}
\usepackage[margin=25mm]{geometry}
\usepackage{fontspec}
\setmainfont{Latin Modern Roman}
\usepackage{amsmath,amssymb,mathtools,unicode-math}
\setmathfont{Latin Modern Math}
\usepackage{enumitem,needspace,xurl,hyperref}
\setlength{\parindent}{0pt}
\setlength{\parskip}{5pt}
\setlength{\emergencystretch}{3em}
\newcommand{\N}{\mathbb{N}}
\newcommand{\Z}{\mathbb{Z}}
\newcommand{\Q}{\mathbb{Q}}
\newcommand{\R}{\mathbb{R}}
\newcommand{\C}{\mathbb{C}}
\newcommand{\F}{\mathbb{F}}
\newcommand{\A}{\mathbb{A}}
\newcommand{\PP}{\mathbb{P}}
\newcommand{\E}{\mathbb{E}}
\newcommand{\eps}{\varepsilon}
\newcommand{\dd}{\,\mathrm{d}}
\begin{document}
\section*{UnsolvedMath ID: 20002894}
\textbf{Source problem code:} catalogue entry
\subsection{Definitions and mathematical statement}
Peter Ebenfelt's Problem 12 asks for hypersurface hypotheses forcing transversality of every noncollapsed holomorphic map. Let \(0\in M\subset\C^{n+1}\), \(0\in M'\subset\C^{N+1}\) be \(C^\infty\) or real-analytic real hypersurfaces, \(n\le N\), and let \(H:(\C^{n+1},0)\to(\C^{N+1},0)\) be holomorphic with \(H(M)\subset M'\). Assume no full neighborhood is mapped into \(M'\). Determine conditions on \(M,M'\) under which every such \(H\) is transverse at 0, meaning
\[
 d_0H(T_0\C^{n+1})+T_0M'=T_0\C^{N+1}
\]
as real tangent spaces. More generally the target tangent spaces are based at \(H(x)\). With real defining functions \(\rho,\rho'\), write \(\rho'\circ H=a\rho\); on \(M\), transversality is equivalent to \(a(x)\ne0\).

The source records equal-dimensional finite-map transversality when either hypersurface is of finite type, and transversality when both are strictly pseudoconvex, under its noncollapse assumptions. A finite holomorphic germ here has isolated fiber over 0, with the usual equal-dimensional interpretation in that statement. These are reported background cases, not a complete criterion.

The supplied verification restores the fractions and signs in the source's example:
\[
 H(z,w)=(z+z^2+\tfrac i2w,\ z-z^2-\tfrac i2w,\ -2zw),
\]
\[
 M=\{\operatorname{Im}w=|z|^2\},\qquad
 M'=\{\operatorname{Im}W=-|Z_1|^2+|Z_2|^2\}.
\]
For \(\rho=\operatorname{Im}w-|z|^2\) and \(\rho'=\operatorname{Im}W+|Z_1|^2-|Z_2|^2\), the aid reports the exact identity
\[
 \rho'\circ H=-2(z+\overline z)\rho.
\]
Thus the map is transverse precisely on \(M\setminus\{\operatorname{Re}z=0\}\), and fails at 0 despite not collapsing a neighborhood into the target. The target is indefinite, so this does not contradict the strictly pseudoconvex case.

\textbf{Quadratic ansatz classified in the aid.} For \(a,b,c\in\C\), set
\[
 H_{a,b,c}(z,w)=(z+az^2+bw,\ z-az^2-bw,\ czw).
\]
The aid reports
\[
 H_{a,b,c}(M)\subset M'
 \quad\Longleftrightarrow\quad
 a=2i\overline b,\quad c=-4i\overline b.
\]
Under these conditions,
\[
 \rho'\circ H_{a,b,c}=(cz+\overline{cz})\rho.
\]
If \(b=0\), the map is \((z,z,0)\) and is collapsed into \(M'\). If \(b\ne0\), it is an immersion at 0 and a finite germ of multiplicity one onto its image, with exact nontransversality locus \(\operatorname{Re}(cz)=0\) on \(M\). Finite here is in the positive-codimension, finite-onto-image sense and does not invoke the source's equal-dimensional result. The original example is \((a,b,c)=(1,i/2,-2)\).

The aid also reports post-workshop criteria involving equal-dimensional generic-full-rank maps, Levi-rank/dimension/Jacobian-rank-stratum inequalities in positive codimension, and a 2025 low-codimension theorem for maps into same-signature hyperquadrics. Their exact inequalities, scope and references are absent; no quantitative criterion is invented from those summaries. The classification of the displayed ansatz is not a classification of all holomorphic maps or a solution of the general hypersurface-condition problem.
\subsection{Short English statement}
The source seeks hypotheses making every noncollapsed hypersurface map transverse. The aid verifies an indefinite-target counterexample and classifies a quadratic family containing it, with an exact real multiplier detecting transversality. Positive-codimension finite immersions can still fail transversality; broader reported criteria need their missing hypotheses and references.
\subsection{Sources}
\begin{itemize}
\item Local UnsolvedMath registry, permanent ID 20002894, catalogue entry; canonical \texttt{aim-several-complex-variables-notes.json}, notes[51].
\item AIM \emph{Complexity of mappings in CR geometry}, Problem 12, proposed by Peter Ebenfelt, \url{https://aimath.org/WWN/complexitycrmap/complexitycrmap.pdf}. The supplied verification reports removing a running header and recovering primes, fractions and tangent subscripts by PDF inspection; this draft does not repeat that check.
\item The verification cross-checks the example against published Example 2.4 of Baouendi--Ebenfelt--Rothschild, without a full title, date or bibliographic entry.
\item Machine-generated research attempt 1, partial-result/partially-solved labels, reported rigor proved and low novelty confidence. It mentions post-workshop and 2025 criteria without full citations. Its ansatz classification is not peer reviewed and requires independent verification.
\end{itemize}
\end{document}
